\lstset{basicstyle=\footnotesize\ttfamily}
\section{Appendix}

This appendix contains material not published in the journal article,
including a design for additional language features not yet
implemented, and example programs.  All the example programs are fully
functional with the current implementation.

\subsection{About the Name}

%% \begin{wrapfigure}{r}{0.25\textwidth}
%%   \begin{center}
%%     \includegraphics[width=0.25\textwidth]{Metropolisposter}
%%   \end{center}
%% \end{wrapfigure}

Babelsberg is a district of Potsdam, in which HPI is located.  It is
also the location of Studio Babelsberg, the historical center of the
German film industry (Fritz Lang's \emph{Metropolis} and Josef von
Sternberg's \emph{The Blue Angel} were filmed there), and still
active, primarily for high-budget feature films.

% \begin{figure}[h]
%   \begin{center}
%   \includegraphics[width=0.25\textwidth]{Metropolisposter}
%   \end{center}
%   \caption{Poster for ``Metropolis''}
% \end{figure}

\subsection{Support for Multiple Cooperating Solvers}
\label{sec:cooperating-solvers}

The current implementation supports multiple solvers, and a protocol
for solvers so that it is convenient to add new ones.  However, for a
given variable only one solver is used.  This is a significant
limitation, since richer sets of constraints can be written and solved
using multiple cooperating solvers.  The Viewpoints note
``Architectures for Cooperating Constraint Solvers''
\cite{borning-coop-solvers-2012} describes one technique for
supporting this, which fits well with \langname.
%% Our more recent
%% work with Bloom as an infrastructure for cooperating solvers takes a
%% more general approach, which will also be more complex to implement.
%% We suggest an evolutionary path for \langname that begins with the
%% former approach, and perhaps later encompasses the Bloom approach.
Briefly, in the ``Architectures for Cooperating Constraint Solvers''
approach, we view the constraints and constrained variables as a
bipartite graph.  For the top level, we partition the constraints into
regions, based on the following requirements:
\begin{itemize}
\item each constraint belongs to exactly one region
\item the different regions can share variables (but not constraints)
\item for each variable shared by two or more regions, that variable
  must be read-only for all but one of the regions
\item the graph of the different top-level regions must be acyclic
\item it isn't possible to divide any region into two regions and
  still maintain the other requirements
\end{itemize}

We can always partition the constraint graph this way.  (In the limit,
there is a single region.)  The partitioning is unique.

Given this partitioning, and a solver for the constraints in each
region, it is easy to solve the constraint system as a whole: it is just
a dataflow problem.  We find an order for solving the regions, so that
region B is solved after region A if B is downstream from A in terms
of the read-only variables.

%% merge any of this material that's still relevant into the next
%% section

%% \subsection{Support for Local Propagation Constraints and Solvers}

%% Local propagation solvers are useful since they are efficient and can
%% be used with constraints on arbitrary kinds of values.  They have
%% important limitations as well --- in particular they don't handle
%% simultaneous equations, or constraints that don't uniquely determine
%% values for the constrained variables, such as inequalities.  We
%% believe that it is important to support them in \toolname, and also
%% that they be part of a suite of cooperating solvers, others of which
%% \emph{can} handle such kinds of constraints.

%% To allow a constraint to be used with a local propagation solver, the
%% programmer should supply a set of \emph{constraint methods}, one for
%% each writable variable (\ie a variable referenced by the constraint
%% that isn't read-only).  Evaluating the method should set the
%% corresponding variable to a value that satisfies the constraint.  The
%% constraint should be functional for each of its writable variables,
%% \ie the value for that variable should be unique given values for
%% the other constrained variables.  If a constraint is read-only in some
%% variable, then it doesn't need a method for that
%% variable.\footnote{Prior work has sought to remove the restriction
%%   that the local propagation constraint be functional in the writable
%%   variables, so that for example one can use inequalities with local
%%   propagation.  Other work has included local propagation constraints
%%   whose methods set more than one variable (\eg SkyBlue).  However,
%%   these solvers are more difficult for the programmer using them to
%%   understand, so we use basic version of local propagation here.  Note
%%   that if there is only one method per constraint, with all of the
%%   other constrained variables being read-only, local propagation
%%   reduces to dataflow constraints.}

\subsection{Initializing and Manipulating Constraints}
\label{sec:initializing-constraints}

This subsection describes the protocol for initializing constraints
and the methods available on constraint objects.

Formally, a constraint is just a predicate that we want to have hold.
In the implementation, it is represented as an instance of
\lstinline{Constraint}.  Minimally, a constraint can be defined just
as a predicate, written as a Ruby block.  The block should meet the
following criteria:

\begin{itemize}

\item Evaluating the block returns a boolean --- the constraint is
  that it evaluate to true.

\item The expression in the block should either be free of side
  effects, or if there are side effects, they should be benign (for
  example, doing caching).  

\item Repeatedly evaluating the block should return the same thing (so
  an expression involving a random number generator wouldn't qualify).

\end{itemize}

There should be other available attributes for
\lstinline{Constraint}.  Constraints that will be satisfied
using a local propagation solver need a set of propagation methods,
soft constraints need a priority, and so forth.  

\subsubsection{The Syntax of the always Primitive }
\label{sec:synt-always-prim}
The syntax for the \lstinline{always}
primitive is the same as for a normal method
call, but is given here to document the interpretation of the
arguments. Note that the \lstinline{SOLVER_SPECIFIC_ARGS} rule only
gives examples, and different solver libraries can accept different
optional arguments here, but should have sane defaults, if none are
given. The definition of \lstinline{RUBY_STATEMENT} is omitted
(\cf{\cite{flanagan2008ruby}}).
\begin{lstlisting}{columns=fixed}
ALWAYS: "always" ["(" ALWAYS_ARGS ")"] CONSTRAINT_BLOCK

ALWAYS_ARGS: ["solver:" RUBY_STATEMENT] [SOLVER_SPECIFIC_ARGS]
            | [SOLVER_SPECIFIC_ARGS]

SOLVER_SPECIFIC_ARGS: ["priority": PRIORITY] [SOLVER_ARGS]
                      | [predicate": BLOCK] [SOLVER_ARGS]
                      | // ...

PRIORITY: "required" | "high | "medium" | "low"

CONSTRAINT_BLOCK: "DO" RUBY_STATEMENT* "END"
                  | "{" RUBY_STATEMENT* "}"
\end{lstlisting}

\subsubsection{Initializing a Constraint}

Here is a list of attributes that can be supplied when initializing a
constraint object.  The predicate is the only required one, and as
described below can be either an argument or a block after the
arguments.  All of the other attributes are specifiable as optional
arguments to \lstinline{always} or \lstinline{once}.

\begin{description}
\item[predicate] A block that evaluates to true or false.  
\item[priority] Defaults to \lstinline{:required}.
\item[weight] Defaults to 1.
\item[error] A block that evaluates to the error in satisfying that
  constraint.  The error value is 0 iff the constraint is satisfied.
  Defaults to a simple 0/1 function.
\item[methods] (optional) A dictionary of local propagation methods, indexed by
  the variable whose value is computed by the corresponding
  method.
\end{description}

The local propagation methods are only needed for local propagation
solvers such as DeltaBlue.  Each method is a block, that if evaluated,
will set the index variable to a value that satisfies the constraint.
The constraint should be functional for the writable variables, \ie
the value for the index variable should be unique given values for the
other constrained variables.  If a constraint is read-only in some
variable, then there isn't a corresponding block for that variable in
the dictionary.

\lstinline{always} can be used to define compound constraints. This
happens, for example, if a predicate method checks two conditions and
conjoins them (as, for example, in a \lstinline{CartesianPoint#==}
method).  All arguments apply to all constraint parts that are created
during the execution of \lstinline{always}. So a \lstinline{:low} priority
compound constraint will be added to the solver as multiple primitive
constraints that are all \lstinline{:low} priority.  If local propagation
methods are provided, evaluating a method must satisfy all of the
component constraints of a compound constraint.  The error for a
compound constraint is likely to be problematic unless it is a simple
0/1 function.

%% TODO: fix the discussion of local propagation methods, to account
%% for the current syntax

%% Here's a trivial example for the $x+y=z$ constraint.  

\subsubsection{Methods for the Class Constraint}
\label{sec:methods-class-constraint}

The \lstinline{Constraint} class provides the following methods in
\toolname.  (Here we are referring to ordinary Ruby methods provided
by the class, not the constraint methods for local propagation.)

\begin{description}
\item[enable] Enable all constraints that were asserted in the block.
  (If the constraint is a compound constraint added
  to the solver as multiple primitive constraints, all of these
  primitive constraints are enabled.)
\item[disable] Disable all constraints that were asserted in the block.
\item[priority] Returns the constraint block priority.
\item[priority=] Setter for the constraint block priority. This
  retracts all constraints and re-enables them using the new priority.
\item[constraint\_block] Returns the block object that was passed to
  \lstinline{always}.
\item[solver\_constraints] Returns an array with the solver-specific
  constraint objects. For Cassowary constraints, this is a list of
  \lstinline{LinearEqualities}; for \zthree it is a list of
  \lstinline{Z3Asts}.
\end{description}

\subsection{Constraint Durations}
\label{sec:constraint-durations}

A constraint may not necessarily be active for the entire life of the
program.  \langname provides the following options:
\lstinline{always}, \lstinline{once}, and
\lstinline{assert during} (for constraints that are enforced
while a block is executing and then removed).


The return value of a call to \lstinline{always} is an object of type
\lstinline{Constraint}. If the constraint has no solver, but the
expression already evaluates to true, a warning is printed; if it is
false, an exception is raised.

Currently, all assertions that we want to continuously check,
but cannot yet solve, are implemented in this way. For example, without
enabling any string solvers in the system, 
\lstinline!always { stringA == stringB }! will return \lstinline{true} 
if indeed the two strings
are equal, and whenever either of the strings is assigned, the block is
re-executed to check if the condition still holds. If it doesn't,
an exception is raised.

This same mechanism is used when complex objects for which we have no
solver occur in the constraint expression (even if a primitive
constraint is created in the block). Consider equality between two
Cartesian points, \lstinline{pt1} and \lstinline{pt2}: the constraint
expression will create two primitive constraints, namely that
\lstinline{pt1.x == pt2.x} and \lstinline{pt1.y == pt2.y}. Whenever an
\lstinline{@x} or a \lstinline{@y} of either \lstinline{p1} or
\lstinline{pt2} is assigned, the constraint solver will just calculate
new values for the other instance variables. However, if
\lstinline{pt1} is assigned a new object, the two primitive
constraints are removed from the solver and the constraint block is
executed again, to generate new primitive constraints that 
are equivalent to \lstinline{pt1.x == pt2.x}.
Because \lstinline{pt1} may now be an instance of
a different class, the equality may
be expressed differently. So whenever a complex object occurs in the
execution of a constraint expression, the entire block is retained
as long as the object is, so that the block can be re-executed later
when the object is assigned a new value.

\subsection{Constraint Priorities}

All arguments to \lstinline{always} (except for the block), are simply
passed on to the \lstinline{enable} method of the solver specific
primitive constraint objects asserted during the execution of the
block. \toolname does not have specific constraint priorities built in
(or weights, local propagation methods, or anything else that a particular
solver may use). DeltaBlue and Cassowary---the two solvers that we
currently provide that support priorities---simply map a set of named
priorities to internal priority values: \lstinline{required},
\lstinline{high}, \lstinline{medium}, and \lstinline{low}. 
We currently require that all the solver libraries use the same
priority names.  (Adding a more general mechanism that allows
arbitrarily many priorities, and different names in different
libraries, is a possible extension.  However, in practice we've found
that the fixed set of priorities has been sufficient, because the priorities
tend to be used in stylized ways, \eg \lstinline{high} for edit
constraints and \lstinline{low} for stays.)

%% COMMENTED OUT SINCE NOW MOSTLY REDUNDANT WITH THE CONCLUSION:
%%
%% \subsection{Extending Other Languages}
%% \label{sec:other-languages}

%% \toolname is an extension to Ruby, and Ruby programs without
%% constraints continue to be legal in it.  The Babelsberg syntax takes
%% advantage of some Ruby quirks to provide a convenient syntax, and the
%% implementation (\autoref{sec:implementation}) relies on Topaz, an
%% experimental \ac{vm} built using the PyPy/RPython toolchain.

%% However, the ideas are applicable to other dynamic object-oriented
%% languages, and future work may include versions in Javascript or
%% Squeak.  In particular, a constraint extension to Javascript would be
%% of particular interest given the potential for integration with
%% Lively.

%% Another potential direction for future work is to evolve this language
%% extension to work with the constraint-based extension to Bloom,
%% initially in Ruby, and later in other languages as well to support
%% cooperating languages with constraints.

\subsection{Example Programs}

This part of the appendix provides a set of example programs in
\toolname that illustrate various features of the language.

\subsubsection{Electrical Circuit Examples}
\label{sec:electrical-circuits}

Here is a complete listing of the classes for building simple
electrical circuits with batteries, resistors, wires, and grounds,
along with some sample circuits using them.  Note in particular the
clean way constraints such as Ohm's Law can be expressed, and the
integration of constraints with control structures for the Kirchhoff's
Law constraint that the currents in a list of connected leads sum to 0
in line 62.

\begin{lstlisting}
class Lead
  attr_reader :voltage, :current
  def initialize
    # initialize voltage and current to 0.0 so that they are of 
    # type Float (constraints may change them later)
    @voltage = 0.0
    @current = 0.0
  end
end

class TwoLeadedObject
  attr_reader :lead1, :lead2
  def initialize
    @lead1 = Lead.new
    @lead2 = Lead.new
    # constrain currents to be equal magnitude and opposite
    always { lead1.current + lead2.current == 0.0 }
  end
end

class Resistor < TwoLeadedObject
  attr_reader :resistance
  def initialize(resistance)
    super()
    @resistance = resistance
    # Ohm's Law constraint.  Note that the resistance is read-only.
    always { lead1.voltage - lead2.voltage == resistance.?*lead1.current }
  end
end

class Battery < TwoLeadedObject
  attr_reader :supply_voltage
  def initialize(supply_voltage)
    super()
    @supply_voltage = supply_voltage
    always { lead2.voltage - lead1.voltage == @supply_voltage.? }
  end
end

class Ground
  attr_reader :lead
  def initialize
    @lead = Lead.new
    # constrain the voltage and current to be 0
    always { lead.voltage == 0.0 }
    always { lead.current == 0.0 }
  end
end

class Wire < TwoLeadedObject
  def initialize
    super()
    always { lead1.voltage == lead2.voltage }
  end
end

def connect(*leads)
  return if leads.empty?
  # all voltages should be equal
  leads[1..-1].each { |a| always { a.voltage == leads[0].voltage } }
  # sum of currents has to be 0
  always { leads.map(\&:current).sum == 0 }
end

def battery_resistor
  # build a circuit with a battery and a resistor, and return the current in one of the leads
  r = Resistor.new(resistance)
  b = Battery.new(voltage)
  connect g.lead, r.lead1, b.lead1
  connect r.lead2, b.lead2
  return r.lead1.current
end

def wheatstone_bridge
  # build a Wheatstone bridge, and return the current through the
  # galvanometer (represented as a wire in this version)
  g = Ground.new
  b = Battery.new(5.0)
  r1 = Resistor.new(100.0)
  r2 = Resistor.new(100.0)
  r3 = Resistor.new(100.0)
  r4 = Resistor.new(100.0)
  w = Wire.new
  connect g.lead, b.lead1, r3.lead2, r4.lead2
  connect b.lead2, r1.lead1, r2.lead1
  connect r1.lead2, r3.lead1, w.lead1
  connect r2.lead2, r4.lead1, w.lead2
  return w.lead1.current
end
\end{lstlisting}

%% \subsubsection{Doubly Linked Lists}

%% This example is adapted from one in reference \cite{lopez-ecoop-1994},
%% and illustrates the use of identity constraints. This example
%% currently does \emph{not} work in \toolname, because identity
%% constraints and conditional statements in constraints are not fully
%% implemented.

%% For comparison, first here is the code for a singly linked list.

%% \begin{lstlisting}
%% class ListCell

%%   attr_accessor :value, :next

%%   def initialize(value,next=nil)
%%     @value = value
%%     @next = next
%%   end

%%   def insert_before(new_cell)
%%     new_cell.next = self
%%     new_cell
%%   end

%%   def insert_after(new_cell)
%%     temp = @next
%%     @next = new_cell
%%     new_cell.next = temp
%%     self
%%   end
%% \end{lstlisting}

%% Here is the code for \lstinline{DoublyLinkedListCell} --- the only
%% difference is the additional instance variable \lstinline{prev} and
%% the constraint in \lstinline{initialize}.

%% \begin{lstlisting}
%% class DoublyLinkedListCell

%%   attr_accessor :value, :next, :prev 
%%   # :prev is only used internally

%%   def initialize(value,next=nil)
%%     @value = value
%%     @next = next
%%     always { @next.prev.equal?(self) if @next }
%%   end

%%   # insert_before and insert_after are the same 
%% end
%% \end{lstlisting}
\subsubsection{Cryptarithmetic}

Here is the well known cryptarithmetic puzzle send+more=money, which \toolname can solve
using \zthree.  (Later, we should also have an efficient finite-domain
solver such as kodkod available.)
\begin{lstlisting}
# initialize each variable to an integer so that the solver knows its type
s,e,n,d,m,o,r,y = [0]*8

# each letter represents a digit between 0 and 9
always { [s,e,n,d,m,o,r,y].ins(0..9) } 

# each letter is mapped to a different digit
always { [s,e,n,d,m,o,r,y].alldifferent? }

# the addition constraint
always {   s*1000 + e*100 + n*10 + d +
           m*1000 + o*100 + r*10 + e ==
 m*10000 + o*1000 + n*100 + e*10 + y }

# the leading digits can't be 0
always { s>0 }
always { m>0 }

puts ("solution: [s,e,n,d,m,o,r,y] = " + [s,e,n,d,m,o,r,y].to_s)
\end{lstlisting}


%%% Local Variables: 
%%% mode: latex
%%% TeX-master: "babelsberg_hpi_technical_report"
%%% End: 
